\section{Transport between saddles and maxima}\label{sec:transport}

We now show that the weighted law obtained from $\rho>2$ is impossible.
The pairing of Proposition~\ref{prop:critical-pairing} compares masses
after changing variables by $G_F$.  Averaging that comparison over large
boxes will give $\EE J_F(i)\leq0$, contradicting
Lemma~\ref{lem:jacobian-moments}.

Throughout this section, suppose $\rho>2$, let
$\beta=\log_2\rho$ and $k=(\beta+3)/4$, and use the probability $P$ from
Lemma~\ref{lem:weighted-law}.  Write
\[
 J_{F,+}=\max(J_F,0),\qquad J_{F,-}=\max(-J_F,0),\qquad
 \omega_F(x+iy)=y^{\beta-1}|Q_F(x+iy)|^2.
\]
Recall that $\mathcal G\subset\KK\times\CC$ is the Borel set of good
pairs and that $R_F(z)$ is the assigned maximum whenever
$J_F(z)>0$ and $(F,G_F(z))\in\mathcal G$.

\subsection{Change of variables and exceptional targets}

For fixed $F$, the open set $\{J_F>0\}$ has a countable cover by open
sets on which $G_F$ is a diffeomorphism onto an open set.  Disjointifying
this cover gives Borel pieces $E_j$ on each of which $G_F$ is injective.
Their images are Borel: each restriction is contained in a
diffeomorphism chart, whose inverse is continuous.  Ordinary
absolute-Jacobian change of variables on these pieces, followed by
Tonelli's theorem \cite[Theorems~2.37 and~2.47]{Folland1999},
therefore gives change of variables with a sum over
preimages.  The same construction applies on $\{J_F<0\}$.
In particular, for every nonnegative Borel function $h$ on $\HH$,
\begin{equation}\label{eq:transport-area}
 \int_{\HH}\omega_F(z)J_{F,\pm}(z)h(z)\dd A(z)
 =
 \int_{\CC}
 \sum_{\substack{G_F(z)=\xi\\ \pm J_F(z)>0}}
 y^{\beta-1}|Q_F(z)|^4h(z)\dd A(\xi).
\end{equation}
Indeed, $|\det DG_F|=|F'|^2|J_F|$, so division by the absolute
Jacobian supplies the second factor $|Q_F|^2$.  All sums and integrals
in \eqref{eq:transport-area} are nonnegative; no global injectivity of
$G_F$ or finite total number of preimages is required.

The root target is $G_F(i)=-i/k$ for every $F$.  Since the exceptional
target set can depend on $F$, the almost-everywhere target assertion
in Lemma~\ref{lem:good-pairs} must still be transferred to this root.

\begin{lemma}[Exceptional targets at the root]\label{lem:root-exception}
 Suppose $\rho>2$, let $P$ be the weighted law of
 Lemma~\ref{lem:weighted-law}, and set
 $k=(\log_2\rho+3)/4$.  Then
 \begin{equation}\label{eq:root-exception}
  \EE\!\left[
   J_{F,+}(i)\,
   \mathbf 1_{\{(F,-i/k)\notin\mathcal G\}}
  \right]=0.
 \end{equation}
\end{lemma}

\begin{proof}
 For $P$-almost every $F$, the complement of
 $\{\xi:(F,\xi)\in\mathcal G\}$ has planar area zero by
 Lemma~\ref{lem:good-pairs}.  Apply \eqref{eq:transport-area} on the
 positive sign with
 \[
 h(z)=\mathbf 1_{B_0}(z)
       \mathbf 1_{\{(F,G_F(z))\notin\mathcal G\}},
 \qquad B_0=\{x+iy:|x|\leq1,\ 1\leq y\leq2\}.
 \]
 Its target integral is supported on that null set, so
 \begin{equation}\label{eq:exception-box}
  \int_{B_0}\omega_F(z)J_{F,+}(z)
   \mathbf 1_{\{(F,G_F(z))\notin\mathcal G\}}\dd A(z)=0
 \end{equation}
 for almost every $F$.

 Define the nonnegative Borel function
 \[
  \Phi(F)=J_{F,+}(i)
          \mathbf 1_{\{(F,-i/k)\notin\mathcal G\}}.
 \]
 The normalized Jacobian and goodness are affine-equivariant, so
 \[
  \Phi(T_zF)=J_{F,+}(z)
            \mathbf 1_{\{(F,G_F(z))\notin\mathcal G\}}.
 \]
 The weighted covariance identity consequently gives
 \[
  \EE[\omega_F(z)\Phi(T_zF)]
   =y^{\beta-1}y^{-\beta}\EE\Phi
   =y^{-1}\EE\Phi.
 \]
 Taking expectations in \eqref{eq:exception-box} and using Tonelli
 yields $0=2\log2\,\EE\Phi$, which proves the claim.
\end{proof}

\subsection{Box comparison and exhaustion}

For each positive integer $N$, let $D_N$ be the Borel set of pairs
$(F,z)$ for which $J_F(z)>0$, $(F,G_F(z))\in\mathcal G$, and the
assigned point $z'=R_F(z)=x'+iy'$ satisfies
\begin{equation}\label{eq:transport-displacement}
 \frac1N\leq\frac{y'}y\leq N,\qquad |x'-x|\leq Ny,
 \qquad z=x+iy.
\end{equation}
The Borel property follows from that of the pairing map.
Its equivariance gives
\begin{equation}\label{eq:transport-equivariance}
 \mathbf 1_{D_N}(F,z)=\mathbf 1_{D_N}(T_zF,i).
\end{equation}
The sets $D_N$ increase with $N$.  Every good-target saddle belongs
to some $D_N$, since its assigned point has finite coordinates and
strictly positive height.

Fix $R>0$ and $Y>1$, and put
\begin{align*}
 B_{R,Y}&=\{x+iy:|x|\leq R,\ 1\leq y\leq Y\},\\
 B'_{R,Y,N}&=\{x+iy:|x|\leq R+NY,\ 1/N\leq y\leq NY\}.
\end{align*}
If $z\in B_{R,Y}$ and $(F,z)\in D_N$, then
$R_F(z)\in B'_{R,Y,N}$ by \eqref{eq:transport-displacement}.
At a preimage $G_F(z)=\xi$, the critical-point identity gives
\begin{equation}\label{eq:transport-target-weight}
 y^{\beta-1}|Q_F(z)|^4
   =\left(\frac{V^F_\xi(z)}k\right)^4,
 \qquad 4(k-1)=\beta-1.
\end{equation}
The pairing is injective on each target fiber, and its assigned
maximum has at least the saddle's $V$-value.  Applying
\eqref{eq:transport-area} to both signs and comparing the target
sums with \eqref{eq:transport-target-weight} therefore proves
\begin{equation}\label{eq:transport-box}
 \int_{B_{R,Y}}\omega_FJ_{F,+}\mathbf 1_{D_N}\dd A
 \leq
 \int_{B'_{R,Y,N}}\omega_FJ_{F,-}\dd A.
\end{equation}
This inequality holds for every $F$; the left side selects only
fibers on which the pairing is defined.
Both source integrands are jointly Borel and nonnegative in $(F,z)$,
so Tonelli's theorem applies when we average this fixed-$F$ inequality.

\begin{proposition}[Transfer bound]\label{prop:transfer-bound}
 The dyadic transfer operator satisfies $\rho\leq2$.
\end{proposition}

\begin{proof}
 Continue under the contrary assumption $\rho>2$.  By weighted
 covariance and \eqref{eq:transport-equivariance},
 \[
 \EE[\omega_F(z)J_{F,+}(z)\mathbf 1_{D_N}(F,z)]
  =y^{-1}\EE[J_{F,+}(i)\mathbf 1_{D_N}(F,i)].
 \]
 Likewise $\EE[\omega_F(z)J_{F,-}(z)]=y^{-1}\EE J_{F,-}(i)$.
 Thus expectation in \eqref{eq:transport-box} yields
 \begin{equation}\label{eq:transport-expected-box}
 \begin{split}
 2R\log Y\,
   \EE[J_{F,+}(i)\mathbf 1_{D_N}(F,i)]
 \leq{}&
 2(R+NY)\log(N^2Y)\,\EE J_{F,-}(i).
 \end{split}
 \end{equation}
 The root quantities are integrable: $q_F(i)$ is uniformly bounded
 on the compact normalized class, and $J_F(i)$ is a fixed quadratic
 expression in this jet.

 Hold $N$ fixed, set $R=Y^2$, and let $Y\to\infty$.  The ratio of the
 right box factor to the left one is
 \[
  \left(1+\frac NY\right)
  \left(1+\frac{2\log N}{\log Y}\right)\longrightarrow1.
 \]
 Dividing \eqref{eq:transport-expected-box} by $2R\log Y$ gives
 \[
  \EE[J_{F,+}(i)\mathbf 1_{D_N}(F,i)]\leq\EE J_{F,-}(i).
 \]
 Now let integer $N\to\infty$.  Monotone convergence exhausts the
 good-target saddles, and Lemma~\ref{lem:root-exception} accounts for
 the remaining positive-J weight.  Hence
 \[
  \EE J_{F,+}(i)\leq\EE J_{F,-}(i),\qquad \EE J_F(i)\leq0.
 \]
 For the same $J_F$ and the same law $P$,
 Lemma~\ref{lem:jacobian-moments} gives
 \[
  k^2\EE J_F(i)
  =k(k-1)-(2k-1)\frac{\beta}{2}
                   +\frac{\beta(\beta+1)}4
  =k(k-1)>0,
 \]
 since $2k-1=(\beta+1)/2$ and $k>1$.  This contradiction proves
 $\rho\leq2$.
\end{proof}
